% BEGIN R28_REGIONAL_CLOSURES
% Inserción después de REG01. No es un documento autónomo.
% Procedencia: U012_biografia_estructural_pi.tex, líneas 175--324.
% Testigos orientados: verificar_cierres_regionales.py; mismo problema,
% representantes distintos de los paseos impresos en U012.
% Emisor SHA-256: 42fdacab582de8e2d2cb7dfafc6204d200aef01249b3dcc37949e39507f066d8.
\subsubsection{Multigraphe de blocs et cinq clôtures régionales minimales}
\label{regcl27:seccion}

Les \(468\) émissions de l'ensemble \(\mathcal U_6\), produites par
l'émetteur APP--TRIT--TPK, déterminent un multigraphe de blocs.
L'identification de ses points d'ancrage régionaux permet de résoudre un problème
de couverture conjointe : parcourir les incidences APP, de propagation,
d'auto-échelle et de chaque région de clôture dans une même marche fermée.
Le cycle local de six raccordements de la section~\ref{sec:rec27-empalmes}
et ces couvertures conservent leurs conditions de parcours respectives.

\paragraph{Graphe de blocs et quotient exact de phase.}
Pour \(U=(U_1,\ldots,U_6)\in\mathcal U_6\), nous définissons
\begin{equation}
 A(U)=(U_1,U_2,U_3),\qquad B(U)=(U_4,U_5,U_6).
 \label{regcl27:mitades}
\end{equation}
Le multigraphe non orienté \(\widetilde G\) a pour sommets les triplets
qui apparaissent comme \(A(U)\) ou \(B(U)\), et une arête étiquetée pour chaque
\(U\), incidente à ces deux triplets. Les \(468\) étiquettes sont retenues :
deux émissions distinctes ne sont pas identifiées parce qu'elles partagent leurs extrémités, ni
parce que l'une échange les moitiés de l'autre.

La phase interne agit sur chaque triplet par rotation cyclique :
\begin{equation}
 (a,b,c)\sim(b,c,a)\sim(c,a,b),\qquad
 \operatorname{can}(a,b,c)
 =\min_{\rm lex}\{(a,b,c),(b,c,a),(c,a,b)\}.
 \label{regcl27:fase}
\end{equation}
Le graphe \(G\) remplace les extrémités de l'arête \(U\) par
\(\operatorname{can}A(U)\) et \(\operatorname{can}B(U)\), tout en conservant
son étiquette. L'équivalence agit exactement par les trois rotations
cycliques de chaque sommet ; chaque émission maintient son identité en tant qu'arête.

L'énumération du même émetteur et la recherche de composantes donnent
\begin{equation}
\begin{array}{c|r|l}
 & |V|&\text{tailles des composantes}\\ \hline
 \widetilde G&96&(28,28,28,4,4,4)\\
 G&32&(28,4)
\end{array}
\label{regcl27:componentes}
\end{equation}
Pour obtenir ces données, on réunit les extrémités distinctes, on ajoute
chaque incidence dans les deux sens et on parcourt en largeur chaque
composante non encore visitée. L'adjacence peut être dédupliquée pour compter les
composantes, mais les étiquettes parallèles sont conservées lors de la recherche de fermetures
avec des arêtes prescrites.

\paragraph{Ancrage APP et numérotation régionale.}
Les trois arêtes communes sont
\begin{equation}
\begin{aligned}
 U_{\rm APP}&=903|850|850|252|109|252,\\
 U_e&=850|963|890|149|252|212,\\
 U_\varphi&=830|943|830|109|252|252.
\end{aligned}
\label{regcl27:anclajes}
\end{equation}
Nous fixons l'extrémité \(v_{\rm APP}=\operatorname{can}A(U_{\rm APP})=850|850|903\). Dans les cinq régions déjà construites, on adopte l'ordre suivant, qui place la région transversale sur la cinquième ligne :
\begin{equation}
\begin{array}{c|l|c|r}
 i&U_\pi^{(i)}&\rho_i&m_i\\ \hline
 1&810|923|870|169|272|272&++&144\\
 2&810|983|810|169|272|272&++&144\\
 3&870|923|810|169|272|272&++&144\\
 4&870|923|810|418|674|521&--&144\\
 5&501|614|498|169|272|272&\perp&432
\end{array}
\label{regcl27:regiones}
\end{equation}
La numérotation distingue les trois régions \(++\), le retour muté
\( -- \) et la transversale \(\perp\). Ainsi,
\(\rho=(++ ,++ ,++ ,--,\perp)\) et
\(m=(144,144,144,144,432)\) dans cet ordre ; la somme reste \(1008\).

\paragraph{Cinq témoins complets avec sens de parcours.}
Dans chaque tableau, la ligne \(j\) signifie parcourir l'arête étiquetée \(e_j\) de \(v_{j-1}\) à \(v_j\). Les extrémités sont des représentants canoniques de phase. Puisque \(G\) est non orienté, la paire \((v_{j-1},v_j)\) fixe le sens de parcours de l'arête étiquetée. Chaque tableau contient exactement huit arêtes et satisfait \(v_0=v_8=v_{\rm APP}\).

La recherche en largeur décrite dans la preuve fournit les
témoins suivants, un pour chaque ensemble de quatre arêtes cibles. La
minimalité concerne la longueur ; elle admet plusieurs marches qui la réalisent.
\paragraph{Testigo \(i=1\).}
\begin{center}
\small
\begin{tabular}{rlll}
\hline
\(j\)&\(v_{j-1}\)&\(e_j\)&\(v_j\)\\
\hline
1&\(850|850|903\)&\(109|252|252|850|903|850\)&\(109|252|252\)\\
2&\(109|252|252\)&\(109|252|252|810|923|870\)&\(810|923|870\)\\
3&\(810|923|870\)&\(810|923|870|169|272|272\)&\(169|272|272\)\\
4&\(169|272|272\)&\(169|272|272|850|963|890\)&\(850|963|890\)\\
5&\(850|963|890\)&\(850|963|890|149|252|212\)&\(149|252|212\)\\
6&\(149|252|212\)&\(149|252|212|830|943|830\)&\(830|830|943\)\\
7&\(830|830|943\)&\(830|943|830|109|252|252\)&\(109|252|252\)\\
8&\(109|252|252\)&\(903|850|850|252|109|252\)&\(850|850|903\)\\
\hline
\end{tabular}
\end{center}

\paragraph{Testigo \(i=2\).}
\begin{center}
\small
\begin{tabular}{rlll}
\hline
\(j\)&\(v_{j-1}\)&\(e_j\)&\(v_j\)\\
\hline
1&\(850|850|903\)&\(169|272|272|850|903|850\)&\(169|272|272\)\\
2&\(169|272|272\)&\(169|272|272|810|983|810\)&\(810|810|983\)\\
3&\(810|810|983\)&\(810|983|810|169|272|272\)&\(169|272|272\)\\
4&\(169|272|272\)&\(169|272|272|850|963|890\)&\(850|963|890\)\\
5&\(850|963|890\)&\(850|963|890|149|252|212\)&\(149|252|212\)\\
6&\(149|252|212\)&\(149|252|212|830|943|830\)&\(830|830|943\)\\
7&\(830|830|943\)&\(830|943|830|109|252|252\)&\(109|252|252\)\\
8&\(109|252|252\)&\(903|850|850|252|109|252\)&\(850|850|903\)\\
\hline
\end{tabular}
\end{center}

\paragraph{Testigo \(i=3\).}
\begin{center}
\small
\begin{tabular}{rlll}
\hline
\(j\)&\(v_{j-1}\)&\(e_j\)&\(v_j\)\\
\hline
1&\(850|850|903\)&\(169|272|272|850|903|850\)&\(169|272|272\)\\
2&\(169|272|272\)&\(169|272|272|850|963|890\)&\(850|963|890\)\\
3&\(850|963|890\)&\(850|963|890|149|252|212\)&\(149|252|212\)\\
4&\(149|252|212\)&\(149|252|212|870|923|810\)&\(810|870|923\)\\
5&\(810|870|923\)&\(870|923|810|169|272|272\)&\(169|272|272\)\\
6&\(169|272|272\)&\(169|272|272|830|943|830\)&\(830|830|943\)\\
7&\(830|830|943\)&\(830|943|830|109|252|252\)&\(109|252|252\)\\
8&\(109|252|252\)&\(903|850|850|252|109|252\)&\(850|850|903\)\\
\hline
\end{tabular}
\end{center}

\paragraph{Testigo \(i=4\).}
\begin{center}
\small
\begin{tabular}{rlll}
\hline
\(j\)&\(v_{j-1}\)&\(e_j\)&\(v_j\)\\
\hline
1&\(850|850|903\)&\(169|272|272|850|903|850\)&\(169|272|272\)\\
2&\(169|272|272\)&\(169|272|272|850|963|890\)&\(850|963|890\)\\
3&\(850|963|890\)&\(850|963|890|149|252|212\)&\(149|252|212\)\\
4&\(149|252|212\)&\(149|252|212|870|923|810\)&\(810|870|923\)\\
5&\(810|870|923\)&\(870|923|810|418|674|521\)&\(418|674|521\)\\
6&\(418|674|521\)&\(418|674|521|830|943|830\)&\(830|830|943\)\\
7&\(830|830|943\)&\(830|943|830|109|252|252\)&\(109|252|252\)\\
8&\(109|252|252\)&\(903|850|850|252|109|252\)&\(850|850|903\)\\
\hline
\end{tabular}
\end{center}

\paragraph{Testigo \(i=5\).}
\begin{center}
\small
\begin{tabular}{rlll}
\hline
\(j\)&\(v_{j-1}\)&\(e_j\)&\(v_j\)\\
\hline
1&\(850|850|903\)&\(109|252|252|850|903|850\)&\(109|252|252\)\\
2&\(109|252|252\)&\(109|252|252|501|614|498\)&\(498|501|614\)\\
3&\(498|501|614\)&\(501|614|498|169|272|272\)&\(169|272|272\)\\
4&\(169|272|272\)&\(169|272|272|850|963|890\)&\(850|963|890\)\\
5&\(850|963|890\)&\(850|963|890|149|252|212\)&\(149|252|212\)\\
6&\(149|252|212\)&\(149|252|212|830|943|830\)&\(830|830|943\)\\
7&\(830|830|943\)&\(830|943|830|109|252|252\)&\(109|252|252\)\\
8&\(109|252|252\)&\(903|850|850|252|109|252\)&\(850|850|903\)\\
\hline
\end{tabular}
\end{center}

\begin{theorem}[Minimalité des cinq couvertures régionales]
\label{regcl27:minimalidad}
Pour chaque \(i\in\{1,\ldots,5\}\), parmi les marches fermées de \(G\)
qui parcourent au moins une fois chaque arête de
\[
 E_i^*=\{U_{\rm APP},U_e,U_\varphi,U_\pi^{(i)}\},
\]
la longueur minimale, comptée avec la multiplicité des parcours, est huit.
Les répétitions de sommets et d'arêtes sont permises.
\end{theorem}

\begin{proof}
Toute fermeture qui parcourt \(U_{\rm APP}\) passe par \(v_{\rm APP}\).
La rotation cyclique de sa liste d'étapes permet de commencer en ce sommet
sans changer ni la longueur ni la couverture. Il suffit donc de résoudre le problème
avec une origine fixe.

L'espace de recherche est l'ensemble fini
\[
 \mathscr S_i=V(G)\times\mathcal P(E_i^*),
 \qquad |\mathscr S_i|\le32\cdot16=512.
\]
L'état initial est \(s_0=(v_{\rm APP},\varnothing)\), et l'état
terminal est \(s_*=(v_{\rm APP},E_i^*)\). Lorsqu'on parcourt une arête \(e\)
incidente à \(v\), vers son autre extrémité \(v'\), on met à jour
\begin{equation}
 (v,M)\longmapsto\bigl(v',M\cup(E_i^*\cap\{e\})\bigr).
 \label{regcl27:transicion-bfs}
\end{equation}
Le masque enregistre des étiquettes d'arêtes, et pas seulement des extrémités ; c'est pourquoi des arêtes parallèles distinctes ne doivent pas être fusionnées.

Toutes les transitions coûtent une unité. Une file FIFO, initiée en
\(s_0\), extrait les états par distance croissante, marque chaque état lors de
sa première visite et garde son prédécesseur et l'étiquette parcourue.
De façon équivalente, si \(\mathcal F_0=\{s_0\}\), les nouvelles couches sont
\[
 \mathcal F_{k+1}
 =\operatorname{Succ}(\mathcal F_k)
   \setminus\bigcup_{j=0}^{k}\mathcal F_j.
\]
Par récurrence, \(\mathcal F_k\) contient exactement les états dont la distance minimale à partir de \(s_0\) est \(k\). L'énumération finie sur l'émetteur indiqué donne, pour chacune des cinq régions,
\[
 s_*\notin\bigcup_{k=0}^{7}\mathcal F_k,\qquad s_*\in\mathcal F_8,
 \qquad (\ell_1,\ell_2,\ell_3,\ell_4,\ell_5)=(8,8,8,8,8).
\]
Les tableaux présentent les cinq bornes supérieures, y compris les quatre
étiquettes requises dans chaque cas. L'exploration exhaustive des
couches précédentes prouve la borne inférieure. Leur combinaison prouve la
minimalité dans le problème déclaré.
\end{proof}

La longueur huit caractérise la couverture conjointe de quatre incidences dans chaque région de clôture. Le quotient retient les étiquettes d'émission et leurs raccordements entre classes de phase ; l'évolution temporelle enrichie conserve en outre orientation et mémoire. Cette description relie les cinq régions aux mêmes ancrages de propagation, d'auto-échelle et d'APP, et complète le cycle local de six raccordements. Le prolongement coinductif maintient ses opérateurs et ses cylindres avec ces données d'orientation et de mémoire.
% END R28_REGIONAL_CLOSURES
